\thispagestyle{empty}

\vspace*{\fill}

\begin{center}
    {\scshape Abstract}

\begin{minipage}{0.9\textwidth}

Many scientific machine learning problems involve sequences of function-valued observations. Examples include physical trajectories, spatial fields, demand curves, spectra, and continuous signals.
Existing operator-learning methods largely focus on deterministic prediction, leaving open questions about probabilistic autoregressive models for recursive generation in function space.
We develop a probabilistic autoregressive transformer for this setting and provide a rigorous framework for recursive generation. We show that the model defines a valid stochastic process which, under suitable conditions, converges to a unique invariant law. Furthermore, we study how well a fixed model performs when the training observations are dependent. Two complementary approaches provide generalisation and sample-complexity bounds under several sampling regimes, including independent trajectories, mixing sequences, and continuous functional dynamics. Overall, the thesis answers three questions: whether the model is well-defined, whether recursive generation is stable, and whether the performance of a fixed model can be assessed reliably from dependent data.
 \end{minipage}
\end{center}



\vspace*{\fill}
